\documentclass[11pt]{article}

% -----------------------------------------------------------------------------
% Packages
% -----------------------------------------------------------------------------
\usepackage[margin=1in]{geometry}
\usepackage[T1]{fontenc}
\usepackage[utf8]{inputenc}
\usepackage{lmodern}
\usepackage{microtype}
\usepackage{amsmath,amssymb,amsfonts,mathtools,bm}
\usepackage{booktabs}
\usepackage{longtable}
\usepackage{array}
\usepackage{enumitem}
\usepackage{xcolor}
\usepackage{hyperref}
\usepackage{listings}

\hypersetup{
  colorlinks=true,
  linkcolor=blue!50!black,
  urlcolor=blue!50!black,
  citecolor=blue!50!black,
  pdftitle={Ghost Oracle Converger Whitepaper},
  pdfauthor={Ghost Oracle Suite}
}

\lstdefinestyle{repo}{
  basicstyle=\ttfamily\small,
  breaklines=true,
  columns=fullflexible,
  frame=single,
  rulecolor=\color{black!20},
  xleftmargin=0.5em,
  xrightmargin=0.5em
}

\setlist[itemize]{leftmargin=1.5em}
\setlist[enumerate]{leftmargin=1.8em}

% -----------------------------------------------------------------------------
% Notation
% -----------------------------------------------------------------------------
\newcommand{\R}{\mathbb{R}}
\newcommand{\E}{\mathbb{E}}
\newcommand{\Var}{\operatorname{Var}}
\newcommand{\softmax}{\operatorname{softmax}}
\newcommand{\rank}{\operatorname{rank}}
\newcommand{\tr}{\operatorname{tr}}
\newcommand{\diag}{\operatorname{diag}}
\newcommand{\argmax}{\operatorname*{arg\,max}}
\newcommand{\argmin}{\operatorname*{arg\,min}}
\newcommand{\KL}{D_{\mathrm{KL}}}
\newcommand{\TV}{D_{\mathrm{TV}}}
\newcommand{\norm}[1]{\left\lVert #1 \right\rVert}
\newcommand{\ip}[2]{\left\langle #1,#2 \right\rangle}
\newcommand{\Gm}{G_M}
\newcommand{\Sm}{S_M}
\newcommand{\Ts}{T_S}
\newcommand{\Iml}{I_M^{\mathrm{local}}}
\newcommand{\Imf}{I_M^{\mathrm{field}}}
\newcommand{\Fm}{F_M}
\newcommand{\Dm}{D_M}
\newcommand{\geo}{\mathrm{geo}}
\newcommand{\qproj}{\mathrm{qproj}}
\newcommand{\gproj}{\mathrm{gproj}}
\newcommand{\Converger}{\mathcal{C}}
\newcommand{\Transformer}{\mathcal{T}}

\title{\textbf{Ghost Oracle Converger Whitepaper}\\
\large Ghost-Channel Operators Adjacent to Transformer Systems}
\author{Ghost Oracle Suite}
\date{Working Architecture Draft}

\begin{document}
\maketitle

\begin{abstract}
Ghost Oracle Suite is being reframed around \emph{ghost-channel operators}:
auxiliary, side-channel-like measurement and projection channels that expose
structure not captured by a primary classical score alone. The long-term
architecture maps these operators into a \emph{Converger}: a system adjacent to,
but independent from, a transformer. A transformer produces token states,
embeddings, attention scores, memory candidates, and residual streams. The
Converger consumes those objects and evaluates them through a family of
ghost-channel operators.

Each operator is implemented through the same three substrate paths:
a classical geometric reference path, a QPU-calibrated projection path, and a
GPU/noiseless projection path. The purpose is not to make isolated backend
claims, but to provide a controlled framework in which every operator is tested
under a shared task, shared schema, shared controls, and shared metrics.
\end{abstract}

\tableofcontents
\newpage

% -----------------------------------------------------------------------------
\section{Design Objective}

The project objective is to build an operator framework that scales across three
levels:
\begin{enumerate}
  \item \textbf{Consumer scale:} local models, retrieval, ranking, memory,
  dimensional analysis, and reproducible benchmark probes.
  \item \textbf{Industrial scale:} real-time control systems, sensor fields,
  calibrated hardware bases, feedback surfaces, and parameter tuning under
  strict safety envelopes.
  \item \textbf{Network scale:} local models treated as nodes, where each node
  can represent a human context while remaining interoperable with a larger
  system.
\end{enumerate}

The guiding principle is that the system should preserve equal footing between
local and global interests. A local node must not be erased by the aggregate, and
the aggregate must not be made incoherent by purely local optimization. The
Converger is intended as an operator layer for measuring, projecting, coupling,
and controlling structure across those interests.

% -----------------------------------------------------------------------------
\section{Transformer and Converger}

Let a transformer be represented as a map
\begin{equation}
  \Transformer_\theta : X \longrightarrow H,
\end{equation}
where \(X = (x_1,\dots,x_n)\) is an input token sequence and
\[
  H = \{h_t \in \R^d\}_{t=1}^{n}
\]
is a sequence of hidden states. The transformer also induces attention scores,
candidate retrieval sets, residual streams, and token-level prediction scores.

The Converger is an adjacent system
\begin{equation}
  \Converger_\phi : \mathcal{R}(H) \longrightarrow \mathcal{Y},
\end{equation}
where \(\mathcal{R}(H)\) denotes records derived from the transformer state:
hidden states, attention matrices, candidate keys, memory records, residual
streams, or external physical records. The Converger does not replace the
transformer. It evaluates structure around the transformer.

A useful mental model is:
\begin{center}
\begin{tabular}{ll}
\toprule
\textbf{Transformer object} & \textbf{Converger role} \\
\midrule
Embeddings and hidden states & metric, dimension, and interaction probes \\
Attention scores & generalized metric and field deformation probes \\
Retrieval candidates & ranking and projection benchmark substrate \\
Residual stream & local interaction substrate \\
Sequence-level fields & stress and field-interaction substrate \\
Token/statistical traces & fractal expansion substrate \\
External error records & syndrome and stress substrate \\
\bottomrule
\end{tabular}
\end{center}

% -----------------------------------------------------------------------------
\section{Ghost-Channel Operator Package}

Every ghost-channel operator \(O\) is packaged with the same three implementation
files:
\begin{lstlisting}[style=repo]
operator/
|-- operator_qpu_generate.py      # generate or submit the QPU base record
|-- operator_gpu_generate.py      # generate the GPU/noiseless base record
`-- operator_benchmark.py         # compare geo / qproj / gproj
\end{lstlisting}

Each benchmark exposes the same three substrate paths:
\begin{align}
  O^{\geo}   &: \text{classical geometric reference}, \\
  O^{\qproj} &: \text{QPU-calibrated projection/base channel}, \\
  O^{\gproj} &: \text{GPU or noiseless projection/base channel}.
\end{align}

The benchmark projector evaluates
\begin{equation}
  \mathcal{B}_O =
  \operatorname{Benchmark}\!\left(
    O^{\geo}, O^{\qproj}, O^{\gproj};
    \mathcal{D}, \mathcal{N}, \mathcal{M}
  \right),
\end{equation}
where \(\mathcal{D}\) is the task dataset, \(\mathcal{N}\) is the family of
controls or nulls, and \(\mathcal{M}\) is the metric table.

The discipline is:
\begin{equation}
  \text{same operator}
  + \text{same input schema}
  + \text{same controls}
  + \text{same metrics}
  + \text{three substrate paths}.
\end{equation}

% -----------------------------------------------------------------------------
\section{Operator Stack}

The current stack contains seven operators:
\begin{center}
\begin{longtable}{lll}
\toprule
\textbf{Symbol} & \textbf{Name} & \textbf{Primary role} \\
\midrule
\(\Gm\) & Generalized Metric channel & bounded similarity and projection \\
\(\Sm\) & Syndrome Metric channel & syndrome-spacetime field \\
\(\Ts\) & Stress channel & stress tensor from syndrome gradients \\
\(\Iml\) & Local Interaction channel & pointwise channel coupling \\
\(\Imf\) & Field Interaction channel & nonlocal rank or flow deformation \\
\(\Fm\) & Fractal Expansion channel & multi-scale expansion surface \\
\(\Dm\) & Dimensional Metric channel & rank, dimension, and projection stability \\
\bottomrule
\end{longtable}
\end{center}

The next sections define each operator mathematically and map it to a Converger
component.

% -----------------------------------------------------------------------------
\section{G\_M: Generalized Metric Channel}

\subsection{Converger Component}

\(\Gm\) maps most directly to the transformer's embedding, attention, and
retrieval components. If a transformer produces query states \(q_i\) and key
states \(k_j\), then \(\Gm\) evaluates a bounded similarity or projection
coordinate between them.

\subsection{Classical Geometric Path}

For normalized vectors \(q,k \in \R^d\), the classical reference can be cosine
similarity:
\begin{equation}
  s_{\cos}(q,k) =
  \frac{\ip{q}{k}}{\norm{q}_2 \norm{k}_2}.
\end{equation}

A bounded projection coordinate can be written dimensionwise as
\begin{align}
  c_q^{(\ell)} &= \tanh\!\left(\frac{q^{(\ell)}}{\sigma_\ell}\right), \\
  c_k^{(\ell)} &= \tanh\!\left(\frac{k^{(\ell)}}{\sigma_\ell}\right), \\
  g_\ell(q,k) &= \sqrt{\frac{1+c_q^{(\ell)}c_k^{(\ell)}}{2}},
\end{align}
and aggregated as
\begin{equation}
  \Gm^{\geo}(q,k) =
  \frac{1}{d}\sum_{\ell=1}^{d} g_\ell(q,k).
\end{equation}

\subsection{QPU and GPU Projection Paths}

Let \(P_Q\) be a projection table derived from a QPU base and \(P_G\) a matched
noiseless or GPU-generated projection table. With binning maps
\[
  b_q = b(c_q^{(\ell)}), \qquad b_k = b(c_k^{(\ell)}),
\]
the projected paths are
\begin{align}
  \Gm^{\qproj}(q,k) &=
  \frac{1}{d}\sum_{\ell=1}^{d} P_Q[b_q,b_k],\\
  \Gm^{\gproj}(q,k) &=
  \frac{1}{d}\sum_{\ell=1}^{d} P_G[b_q,b_k].
\end{align}

\subsection{Benchmark Use}

The benchmark projector compares retrieval or ranking metrics such as top-1
accuracy, recall at \(K\), mean reciprocal rank, rank deformation, spike
sensitivity, and runtime.

% -----------------------------------------------------------------------------
\section{S\_M: Syndrome Metric Channel}

\subsection{Converger Component}

\(\Sm\) maps to field-like records: sequence fields, error records, repeated
measurements, and time-indexed diagnostics. In transformer-adjacent workloads,
it can be applied to token-time fields, memory-update fields, or external sensor
records.

\subsection{Syndrome Field}

Let
\[
  S \in \{0,1\}^{N \times T \times E}
\]
denote a syndrome-spacetime record with shots \(N\), rounds \(T\), and edges
\(E\). Let final data bits be
\[
  D \in \{0,1\}^{N \times V}.
\]
For adjacent data bits, define terminal edge parity
\begin{equation}
  E_i = D_i \oplus D_{i+1}.
\end{equation}

A simple syndrome metric can measure agreement between terminal edge parity and
syndrome history:
\begin{equation}
  \Sm(S,D)
  =
  \E_{n,t,i}
  \left[
    \mathbf{1}\{S_{n,t,i}=E_{n,i}\}
  \right].
\end{equation}

More generally, \(\Sm\) is a field-valued operator:
\begin{equation}
  \Sm : (S,D) \longrightarrow F_S(t,i),
\end{equation}
where \(F_S\) is a syndrome-spacetime field.

\subsection{Substrate Paths}

\begin{align}
  \Sm^{\geo}   &: \text{classical or synthetic field reference},\\
  \Sm^{\qproj} &: \text{field derived from QPU syndrome records},\\
  \Sm^{\gproj} &: \text{matched GPU/noiseless syndrome field}.
\end{align}

\subsection{Benchmark Use}

Metrics include real-versus-control classification, control-source
classification, distance prediction, field correlation, terminal agreement, and
control collapse under shuffling or reversal.

% -----------------------------------------------------------------------------
\section{T\_S: Stress Channel}

\subsection{Converger Component}

\(\Ts\) maps to stress, curvature, anisotropy, and local gradient behavior in a
field. It is derived from \(\Sm\) but treated as its own operator.

\subsection{Gradient Definitions}

Given a syndrome field \(S(t,i)\), define temporal and spatial finite
differences over \(\mathbb{F}_2\):
\begin{align}
  \Delta_t S(t,i) &= S(t+1,i) \oplus S(t,i),\\
  \Delta_x S(t,i) &= S(t,i+1) \oplus S(t,i).
\end{align}

The stress tensor is
\begin{equation}
  \Ts =
  \begin{bmatrix}
    T_{tt} & T_{tx}\\
    T_{xt} & T_{xx}
  \end{bmatrix},
\end{equation}
where
\begin{align}
  T_{tt} &= \E[(\Delta_t S)^2],\\
  T_{xx} &= \E[(\Delta_x S)^2],\\
  T_{tx} &= T_{xt} = \E[\Delta_t S \, \Delta_x S].
\end{align}

Useful derived quantities are
\begin{align}
  \operatorname{trace}(\Ts) &= T_{tt}+T_{xx},\\
  \operatorname{anisotropy}(\Ts) &= T_{tt}-T_{xx},\\
  \operatorname{coupling}(\Ts) &= T_{tx}.
\end{align}

\subsection{Substrate Paths}

\begin{align}
  \Ts^{\geo}   &: \text{stress of a synthetic or analytical reference field},\\
  \Ts^{\qproj} &: \text{stress of a QPU-derived syndrome field},\\
  \Ts^{\gproj} &: \text{stress of a matched GPU/noiseless field}.
\end{align}

\subsection{Benchmark Use}

Metrics include trace separation, anisotropy, coupling, local trace maps,
classification against controls, and stability under field-preserving nulls.

% -----------------------------------------------------------------------------
\section{I\_M Local: Local Interaction Channel}

\subsection{Converger Component}

\(\Iml\) maps to local interaction between two coordinates, such as a
complexity-like projection coordinate and a countercomplexity-like field
coordinate. In transformer terms, this can be applied to token positions,
candidate ranks, residual features, or local memory rows.

\subsection{Definition}

Let \(Q_i\) be a local generalized metric coordinate and \(F_i\) a local field
coordinate. Define rowwise or local standardization
\begin{equation}
  z(a_i) = \frac{a_i-\mu_a}{\sigma_a+\varepsilon}.
\end{equation}

The local interaction is
\begin{equation}
  {\Iml}_i(Q,F) = z(Q_i)\,z(F_i).
\end{equation}

A score with local interaction can be written
\begin{equation}
  E_i(\lambda) = z(Q_i) + \lambda {\Iml}_i(Q,F).
\end{equation}

\subsection{Substrate Paths}

\begin{align}
  {\Iml}^{\geo}   &: \text{local coupling using classical reference channels},\\
  {\Iml}^{\qproj} &: \text{local coupling using QPU-calibrated channels},\\
  {\Iml}^{\gproj} &: \text{local coupling using GPU/noiseless channels}.
\end{align}

\subsection{Benchmark Use}

Metrics include lift over either channel alone, collapse under shuffled fields,
agreement between QPU and GPU bases, stability under \(\lambda\)-sweeps, and
local rank changes.

% -----------------------------------------------------------------------------
\section{I\_M Field: Field Interaction Channel}

\subsection{Converger Component}

\(\Imf\) maps to nonlocal deformation: rank flow, trajectory changes, sequence
field perturbations, or retrieval-order deformation.

\subsection{Field Deformation}

Let \(s_i\) be a base score along an ordered candidate or field axis. Define a
roughness term
\begin{equation}
  r_i = |s_i-s_{i-1}| + |s_{i+1}-s_i|.
\end{equation}
The field-deformed score is
\begin{equation}
  {\Imf}_i(s;\lambda) = s_i + \lambda z(r_i).
\end{equation}

For a ranking \(\pi\), rank deformation can be measured as
\begin{equation}
  \Delta_{\rank} =
  \frac{1}{m}\sum_{i=1}^{m}
  \left|
    \rank_{\pi_{\mathrm{base}}}(i)
    -
    \rank_{\pi_{\mathrm{field}}}(i)
  \right|.
\end{equation}

\subsection{Substrate Paths}

\begin{align}
  {\Imf}^{\geo}   &: \text{field deformation over classical reference scores},\\
  {\Imf}^{\qproj} &: \text{field deformation using QPU-calibrated scores},\\
  {\Imf}^{\gproj} &: \text{field deformation using GPU/noiseless scores}.
\end{align}

\subsection{Benchmark Use}

Metrics include rank deformation, retrieval lift, trajectory improvement,
over-steering thresholds, \(\lambda\)-stability, and collapse under field
controls.

% -----------------------------------------------------------------------------
\section{F\_M: Fractal Expansion Channel}

\subsection{Converger Component}

\(\Fm\) maps to scale, digit, partition, and expansion behavior in records. It
is the corrected framing of earlier Benford-style analysis: the object is not a
single recursive Benford score, but a multi-scale expansion surface.

\subsection{First-Digit Distribution}

For a base \(b\geq 2\), the Benford first-digit distribution is
\begin{equation}
  P_b(d) = \log_b\!\left(1+\frac{1}{d}\right),
  \qquad d=1,\dots,b-1.
\end{equation}

Given a record \(x\), let \(\widehat{P}_{b,s,c}(x)\) denote the empirical
first-digit or partition distribution at base \(b\), scale \(s\), and chunk
depth \(c\). The fractal signature is
\begin{equation}
  \Phi_F(x) =
  \left\{
    \widehat{P}_{b,s,c}(x)
    :
    b\in\mathcal{B},\;
    s\in\mathcal{S},\;
    c\in\mathcal{C}
  \right\}.
\end{equation}

A distance from control signatures is
\begin{equation}
  d_F(x,\nu) =
  \norm{\Phi_F(x)-\Phi_F(\nu)}_2,
\end{equation}
where \(\nu\) is a matched null or control.

A standardized expansion score is
\begin{equation}
  z_F(x) =
  \frac{
    d_F(x,\bar{\nu}) - \mu_{\nu}
  }{
    \sigma_{\nu}+\varepsilon
  }.
\end{equation}

\subsection{Substrate Paths}

\begin{align}
  \Fm^{\geo}   &: \text{fractal expansion of classical reference records},\\
  \Fm^{\qproj} &: \text{fractal expansion of QPU-derived base records},\\
  \Fm^{\gproj} &: \text{fractal expansion of GPU/noiseless base records}.
\end{align}

\subsection{Benchmark Use}

Metrics include expansion distance, expansion \(z\)-score, base stability,
scale stability, real-versus-control AUC, nearest surviving control, and
\(p\)-adic artifact diagnostics.

% -----------------------------------------------------------------------------
\section{D\_M: Dimensional Metric Channel}

\subsection{Converger Component}

\(\Dm\) maps to dimensional structure: intrinsic dimension, rank collapse,
spectral entropy, projection stability, slicing, lifting, and embedding
compression.

\subsection{Spectral Signature}

Let \(X \in \R^{n\times d}\) be a matrixized record. Define the covariance or
Gram matrix
\begin{equation}
  C = \frac{1}{n}X^\top X.
\end{equation}
Let eigenvalues be
\[
  \lambda_1 \geq \lambda_2 \geq \dots \geq \lambda_d \geq 0.
\]
Normalize them as
\begin{equation}
  p_i = \frac{\lambda_i}{\sum_{j=1}^{d}\lambda_j + \varepsilon}.
\end{equation}

The spectral entropy and effective rank are
\begin{align}
  H_\lambda &= -\sum_{i=1}^{d} p_i \log(p_i+\varepsilon),\\
  r_{\mathrm{eff}} &= \exp(H_\lambda).
\end{align}

The participation ratio is
\begin{equation}
  r_{\mathrm{PR}} =
  \frac{\left(\sum_i \lambda_i\right)^2}
       {\sum_i \lambda_i^2+\varepsilon}.
\end{equation}

The stable rank is
\begin{equation}
  r_{\mathrm{stable}} =
  \frac{\norm{X}_F^2}{\norm{X}_2^2+\varepsilon}.
\end{equation}

Projection stability at target dimension \(k\) can be estimated by a random
projection \(R_k \in \R^{d\times k}\):
\begin{equation}
  \operatorname{stab}_k(X)
  =
  \operatorname{corr}
  \left(
    \operatorname{dist}(X),
    \operatorname{dist}(XR_k)
  \right).
\end{equation}

\subsection{Substrate Paths}

\begin{align}
  \Dm^{\geo}   &: \text{dimension signature of classical reference matrices},\\
  \Dm^{\qproj} &: \text{dimension signature of QPU-derived base matrices},\\
  \Dm^{\gproj} &: \text{dimension signature of GPU/noiseless base matrices}.
\end{align}

\subsection{Benchmark Use}

Metrics include effective rank, spectral entropy, participation ratio, stable
rank, top-eigenvalue fraction, rank thresholds, and random projection stability.

% -----------------------------------------------------------------------------
\section{Localized Training Objective}

The Converger can be used as a localized training or adaptation layer without
requiring that the transformer itself be replaced. Let \(L_{\mathrm{task}}\) be
the original task loss of a model, and let \(L_O\) be an operator loss or
regularizer for ghost-channel operator \(O\). A generic localized objective is
\begin{equation}
  L_{\mathrm{total}}
  =
  L_{\mathrm{task}}
  +
  \sum_{O\in\mathcal{O}}
  \alpha_O L_O
  +
  \beta L_{\mathrm{control}},
\end{equation}
where
\[
  \mathcal{O}
  =
  \{\Gm,\Sm,\Ts,\Iml,\Imf,\Fm,\Dm\}.
\]

A channel-specific operator loss can compare substrate variants:
\begin{equation}
  L_O =
  w_{\geo} L(O^{\geo})
  +
  w_{\qproj} L(O^{\qproj})
  +
  w_{\gproj} L(O^{\gproj}).
\end{equation}

For benchmark discipline, controls should be part of the objective:
\begin{equation}
  L_{\mathrm{control}}
  =
  \max\{0,\; m - \Delta(O_{\mathrm{real}},O_{\mathrm{control}})\},
\end{equation}
where \(m\) is a margin and \(\Delta\) measures separation between real and
control channels.

% -----------------------------------------------------------------------------
\section{Node-Level Convergence}

At network scale, let each local model or human-context model be a node
\(v_i\). Each node has local state \(h_i\), local objectives \(u_i\), and a
Converger signature \(c_i\):
\begin{equation}
  c_i =
  \Converger_\phi(h_i).
\end{equation}

A global coordination objective can be written as
\begin{equation}
  J =
  \sum_i a_i J_{\mathrm{local}}(v_i)
  +
  \gamma J_{\mathrm{global}}(V)
  -
  \eta J_{\mathrm{erasure}}(V),
\end{equation}
where \(J_{\mathrm{erasure}}\) penalizes coordination that destroys local
context or equal footing.

The design goal is not to maximize the aggregate by ignoring nodes, nor to
maximize local objectives by destroying coordination. The goal is convergence
under preserved local representation.

% -----------------------------------------------------------------------------
\section{Hardware Scaling}

The same operator package should scale across hardware levels:
\begin{center}
\begin{tabular}{lll}
\toprule
\textbf{Scale} & \textbf{Hardware} & \textbf{Converger use} \\
\midrule
Consumer & CPU / consumer GPU & local model probes and benchmarks \\
Research & workstation GPU / QPU access & calibrated projection bases \\
Industrial & sensor and actuator control hardware & real-time field tuning \\
Network & model-serving infrastructure & node administration and coordination \\
\bottomrule
\end{tabular}
\end{center}

The key requirement is that the operator schema remains stable:
\[
  \texttt{qpu\_generate}
  \quad
  \texttt{gpu\_generate}
  \quad
  \texttt{benchmark}.
\]

% -----------------------------------------------------------------------------
\section{Claims Discipline}

A ghost-channel claim should not be accepted because a metric is exciting. It
must survive controls. The required pattern is:
\begin{enumerate}
  \item freeze the record;
  \item construct matched controls;
  \item scramble or destroy the proposed channel;
  \item preserve the primary task where possible;
  \item compare \(\geo\), \(\qproj\), and \(\gproj\) under the same metrics;
  \item report both positive results and failure modes.
\end{enumerate}

A channel is load-bearing only if destroying that channel destroys the effect
while preserving the relevant primary structure.

% -----------------------------------------------------------------------------
\section{The Pact}

Ghost Oracle Suite is built around a simple pact:
\begin{quote}
Nothing developed here is meant to compete with intelligence. It is meant to
enhance it.
\end{quote}

That includes human intelligence, local AI systems, large AI systems,
scientific tools, industrial systems, and the people represented by them. The
long-term goal is not to build machinery that replaces the individual with the
aggregate, or the aggregate with the individual. The goal is to build operator
frameworks that help both remain visible.

The design ethic is:
\begin{align}
  \text{design ethic}
  &=
  \text{individual dignity}
  + \text{local context}
  + \text{collective coordination}
  \nonumber\\
  &\quad
  + \text{scientific controls}
  + \text{hardware scalability}
  + \text{open adoption}.
\end{align}

% -----------------------------------------------------------------------------
\section{Conclusion}

The Converger is a transformer-adjacent, hardware-scalable operator layer. It
does not replace the transformer. It evaluates structure around the transformer
through seven ghost-channel operators:
\[
  \Gm,\quad
  \Sm,\quad
  \Ts,\quad
  \Iml,\quad
  \Imf,\quad
  \Fm,\quad
  \Dm.
\]

Each operator is evaluated through three substrate paths:
\[
  O^{\geo},\quad O^{\qproj},\quad O^{\gproj}.
\]

Each operator package contains the same implementation pattern:
\[
  \texttt{qpu\_generate},\quad
  \texttt{gpu\_generate},\quad
  \texttt{benchmark}.
\]

The result is a unified system adjacent to and independent from a transformer:
a Converger that can operate on local consumer hardware, scale to physical
control systems, and eventually support networks of local model nodes without
erasing the humans or contexts those nodes represent.

\end{document}
